\Article{22}
Since the method set out in this memoir finds particularly frequent and convenient application in calculations pertaining to higher geodesy, we hope that readers will welcome an illustration of its rules by some examples drawn from this field.

The conditional equations among the angles of a system of triangles are to be sought chiefly from three sources.

\par\noindent\textbf{I.} The sum of the horizontal angles that complete a full turn of the horizon around the same vertex must equal four right angles.

\par\noindent\textbf{II.} The sum of the three angles in any triangle is equal to a given quantity, since, whenever a triangle lies on a curved surface, the excess of this sum over two right angles can be computed so accurately that it may be treated as absolutely exact.

\par\noindent\textbf{III.} The third source is the ratio of the sides in triangles forming a closed chain. Namely, if a series of triangles is connected so that the second triangle has one side \(a\) in common with the first and another, \(b\), with the third; likewise, if the fourth triangle has a side \(c\) in common with the third and a side \(d\) in common with the fifth, and so on up to the last triangle, whose side in common with the preceding triangle is \(k\), and with the first triangle again is \(l\), then the values of the ratios \(\frac al,\allowbreak\frac ba,\allowbreak\frac cb,\allowbreak\frac dc\ .\ .\ .\ .\ \frac lk\) will be known, respectively, from pairs of angles in successive triangles opposite their common sides, by known methods. Since the product of those fractions must be \(=1\), a conditional equation between the sines of those angles will result (with these angles respectively diminished by the third part of the spherical or spheroidal excess, if the triangles lie on a curved surface).

In more complicated systems of triangles, however, it often happens that more conditional equations of the second and third kinds present themselves than may be retained, since some of them are already contained in the others. The rarer case is that equations similar to the conditional equations of the second kind, concerning figures with more sides, must be added, namely only then; when \SourceLeaf{087}
polygons are formed that have not been divided into triangles by measurement. But these matters, which are too remote from the present purpose, we shall discuss more fully on another occasion. Still, we cannot omit a caution: if our theory is to be applied purely and rigorously, it presupposes that the quantities denoted by \(v,v',v''\) etc. are either actually observed directly or derived from observations in such a way that they remain independent of one another, or at least may be regarded as such. In ordinary practice the angles of the triangles themselves are observed, and so may be taken as \(v,v',v''\) etc. But we must remember that if the system also contains triangles whose angles have not been observed directly, but arise as sums or differences of angles actually observed, those angles must not be counted among the observed quantities; their composite form must instead be retained in the calculations. Matters are different in an observing method like that followed by the distinguished \textsc{Struve} (Astronomische Nachrichten II, p. 431), in which the directions of the individual sides issuing from the same vertex are obtained by comparison with one and the same arbitrary direction. In that case these very angles are to be taken as \(v,v',v''\) etc.; thus all the angles of the triangles appear in the form of differences, and the conditional equations of the first kind, which are automatically satisfied by the nature of the case, cease to be needed. The observing method that I myself followed in the triangulation completed in earlier years differs from both the first and the second method, but in its effect it can be equated with the latter, so that at each station the directions of the sides issuing from it, counted from a starting direction chosen, as it were, arbitrarily, must be taken as the quantities \(v,v',v''\) etc. We shall now work through two examples, one belonging to the first method and the other to the second.
